\section{Introduction}
\label{sec:introduction}

Magnetometer records combine the field of interest with environmental interference, instrument response, and measurement noise. A change in one record can therefore admit several explanations. Measurements from neighboring nodes can provide context, but redundancy alone does not make every environmental and instrumental cause identifiable. Multisensor fusion provides methods for combining such information~\cite{hall1997fusion}. Quantum sensing concerns a different layer: the use of quantum systems to measure physical quantities~\cite{degen2017sensing}. Here, the quantum computation acts on classical measurement features after readout. It does not alter the sensing interaction or establish an improvement in hardware sensitivity.

Quantum kernel methods provide one route for processing these features. A parameterized circuit maps an input vector to a quantum state, and state overlaps define a similarity function for a classical kernel predictor~\cite{schuld2019feature,havlicek2019supervised}. Training the circuit changes that similarity function. This creates two separate questions: whether optimization changes the representation, and whether the resulting representation improves prediction on unseen observations. A large state space answers neither question. The predictive value of a quantum kernel depends on the data and target, as well as on the available classical alternatives~\cite{huang2021power}. Benchmarking studies have reported task-dependent outcomes rather than a general benefit from quantum kernel training~\cite{alvarez2025benchmarking,schnabel2025scrutiny}.

AQSE is a hybrid quantum--classical framework for studying these questions in sensor-data processing. Its network demonstrator connects simulated observations from one to eight magnetometers to causal, eight-dimensional local or network feature profiles. An eight-qubit variational quantum circuit (VQC), with 16 trainable parameters, defines a trainable quantum kernel (TQK). The demonstrator converts kernel similarities to a fixed-dimensional classical representation using a regularized Nystr\"om map, termed AFSE within AQSE, and passes that representation to a multilayer perceptron (MLP). AFSE is an application of a classical kernel approximation, not a separate quantum algorithm~\cite{williams2001nystrom}. Training and inference are distinct operations. Quantum natural gradient (QNG) updates the circuit parameters during an explicit training job~\cite{stokes2020qng}; inference uses frozen parameters and preprocessing. Versioned model bundles link each output to its feature profile, scaler, circuit parameters, landmarks, and classifier~\cite{aqse2026canonical}.

The framework and the statistical benchmark have different scopes. The replicated experiment reported here isolates the TQK component on a controlled, single-magnetometer task: discrimination between nominal and elevated white-noise regimes under harmonic excitation. It uses an eight-feature harmonic representation and a support vector classifier (SVC) operating directly on the quantum kernel. It does not evaluate the complete network State8--AFSE--MLP path. The network architecture establishes the application context; the replicated benchmark tests a bounded component-level hypothesis. All quantum computations use exact statevector simulation. No quantum processing unit (QPU) or physical sensor deployment is evaluated.

The primary question is whether validation-selected TQK models outperform a radial basis function SVC (RBF-SVC) on held-out balanced accuracy within this simulated task. The protocol was versioned in the repository before execution. It specifies 30 independently generated primary replicas, each with 72 training, 24 validation, and 24 test episodes. Each episode contributes one fixed analysis window. Circuit optimization uses a balanced subset of 32 training examples; the final kernel SVC uses all 72 training examples. Preprocessing is fitted on training data, and checkpoint and hyperparameter selection uses validation data. Test evaluation follows the selection freeze. The comparator set includes RBF-SVC, a compact MLP, random Fourier features (RFF) with a linear SVC, and gradient boosting. Two additional control sets contain 12 runs each~\cite{aqse2026canonical}.

This paper makes three contributions. First, it describes a versioned sensor-processing framework that separates observable measurements, simulator-only causes, learned representations, and predictions. Second, it reports a repository-preregistered comparison of TQK learning and classical alternatives under a common episode-level data split. Third, it examines kernel spectra across optimization checkpoints and reports the controls alongside predictive results. These diagnostics describe representation and selection behavior; they are not substitutes for held-out evaluation.

The mean test balanced accuracy was 0.9528 for TQK and 0.9569 for RBF-SVC. Their paired difference was $-0.0042$, with a 95\% bootstrap interval of $[-0.0222,\,0.0125]$ and a two-sided Wilcoxon $p=0.933$. The preregistered predictive-advantage criterion was not met~\cite{aqse2026canonical}. The training-label permutation control also requires qualification: validation labels remained available for model selection. It is therefore a partial training-signal ablation, not a complete removal of supervised information. We interpret the results within these boundaries and distinguish changes in kernel geometry from evidence of predictive or computational advantage.
